\documentclass[12pt]{article}
\usepackage[utf8]{inputenc}
\usepackage[T1]{fontenc}
\usepackage[margin=1in]{geometry}
\usepackage{setspace}
\usepackage{xcolor}
\usepackage[colorlinks=true, urlcolor=blue!60!black, linkcolor=blue!60!black]{hyperref}
\setlength{\parindent}{2em}
\setlength{\parskip}{1.25em}
\renewcommand{\baselinestretch}{1.0}
\begin{document}
\begin{center}
{\LARGE \textit{Virginia State Board of Pharmacy v. Virginia Citizens Consumer Council}, 425 U.S. 748 (1976)\par}
\vspace{0.8em}
{\large Case Brief --- Four-Page Edition\par}
\end{center}
\vspace{1em}

\section{Memory Jogger}

Striking down Virginia's ban on advertising prescription-drug prices, the Court held 7--1 that purely commercial speech is protected by the First Amendment and that consumers have a reciprocal right to receive truthful information about lawful products (\href{https://supreme.justia.com/cases/federal/us/425/748/}{425 U.S.\ 748}).

\section{Detailed Case Facts}

Virginia law made it ``unprofessional conduct,'' punishable by fines or license suspension or revocation, for a pharmacist to advertise any price, discount, or credit terms for prescription drugs. Because only pharmacists could dispense such drugs, the statute effectively banned all prescription-drug price advertising, even though about 95\% of prescriptions were filled with prepackaged, manufacturer-prepared drugs.

A 1969 federal court had upheld the statute against due process and equal protection claims by a retailer and a pharmacist. This suit was brought by \emph{consumers}: a Virginia resident who needed daily prescription medication, the Virginia Citizens Consumer Council, and the Virginia State AFL-CIO. The parties stipulated that prices varied dramatically within the same localities---40 Achromycin tablets in Richmond ranged from \$2.59 to \$6.00, and tetracycline in Newport News--Hampton from \$1.20 to \$9.00---and that some pharmacists would advertise if allowed.

\section{Procedural History}

A three-judge district court in the Eastern District of Virginia declared the statute void under the First and Fourteenth Amendments and enjoined it. The Board of Pharmacy appealed directly. The Court heard argument on November 11, 1975, and affirmed on May 24, 1976.

\section{Judicial Votes}

Affirmed, \textbf{7--1}.
\begin{itemize}
\item \textbf{Majority:} Justice Blackmun, joined by Chief Justice Burger and Justices Brennan, Stewart, White, Marshall, and Powell.
\item \textbf{Concurring:} Chief Justice Burger; Justice Stewart.
\item \textbf{Dissenting:} Justice Rehnquist.
\item \textbf{Not participating:} Justice Stevens.
\end{itemize}

\section{Holding}

Where a willing speaker exists, First Amendment protection extends to the communication, its source, and its recipients, so consumers have a reciprocal right to receive truthful advertising. Purely commercial speech is not wholly outside the First Amendment, and a State may not completely suppress truthful information about entirely lawful activity out of fear of its effect on speakers and listeners.

\section{Analysis of Opinions}

\textbf{Majority (Blackmun).} Blackmun dismantled the commercial-speech exception dating to \textit{Valentine v.\ Chrestensen} (1942), extending \textit{Bigelow v.\ Virginia}. \emph{Listeners:} under \textit{Lamont}, \textit{Kleindienst v.\ Mandel}, and \textit{Procunier v.\ Martinez}, the First Amendment protects the listener's reciprocal right. \emph{Value:} speech that does ``no more than propose a commercial transaction'' still matters, as labor-dispute speech did in \textit{Thornhill}. The consumer's interest is acute for ``the poor, the sick, and particularly the aged,'' and society has an interest in intelligent private resource allocation, which also informs debate over economic regulation. \emph{Paternalism:} Virginia could regulate pharmacists directly, but its ban worked only by keeping citizens ignorant on the ``highly paternalistic'' assumption that they would choose badly. \emph{Limits:} the State may still impose time, place, and manner rules and ban false, misleading, or unlawful ads. Footnote~24 noted commercial speech's greater verifiability and durability, and footnote~25 reserved lawyer and doctor advertising.

\textbf{Burger (concurring).} Pharmacists filling prepackaged prescriptions are largely retail packagers; advertising by the ``traditional learned professions of medicine or law'' poses greater risks.

\textbf{Stewart (concurring).} Because commercial claims are verifiable, banning false ads does not chill accurate information, unlike the ``breathing space'' ideological speech needs under \textit{New York Times v.\ Sullivan}.

\textbf{Rehnquist (dissenting).} Consumers could still ask prices themselves; the Court revived \textit{Lochner} under a First Amendment guise; and its logic would reach lawyers, liquor, cigarettes, and ads urging patients to ``Insist that your physician prescribe Demerol.''

\textbf{Later developments.} \textit{Bates v.\ State Bar of Arizona} (1977) extended the rule to lawyers; \textit{Central Hudson} (1980) set a four-part test; \textit{44 Liquormart} and \textit{Lorillard} hardened anti-paternalism; and \textit{Sorrell v.\ IMS Health} (2011) applied ``heightened judicial scrutiny'' to content- and speaker-based burdens.

\section{Examples: Future Applications}

\begin{enumerate}
\item \textbf{Same-Side (Invalid): Optometrist Price Ads.} A ban on advertising eyeglass prices to protect professional standards keeps the public in ignorance, even though \textit{Williamson v.\ Lee Optical Co.} once upheld a similar rule.
\item \textbf{Same-Side (Invalid): Lawyers' Flat Fees.} Disciplining lawyers for advertising flat fees for routine services fails under \textit{Bates}.
\item \textbf{Opposite-Side (Valid): Solicitation.} Pressuring a hospitalized accident victim to sign on the spot may be banned under \textit{Ohralik} (1978).
\item \textbf{Opposite-Side (Valid): Ads for Illegal Sales.} Advertising prescription drugs without a prescription proposes an unlawful transaction and is unprotected.
\item \textbf{Fence-Sitter: Sportsbook Promotions.} A ban on advertising signup bonuses, with terms still shown in-app, restricts truthful information to reduce demand, but targets a vice with real harms and leaves a channel open; \textit{Greater New Orleans Broadcasting} requires evidence of fit.
\end{enumerate}

\section{Critique}

\textbf{Free-market constitutionalism versus self-government.} Blackmun linked advertising to Alexander Meiklejohn's self-government theory by arguing that informed economic choices shape public regulatory choices. Critics respond that this collapses political participation into consumer preference and constitutionalizes laissez-faire assumptions.

\textbf{Rehnquist's prescience.} His Demerol and Seconal hypotheticals anticipated today's ubiquitous direct-to-consumer drug advertising and the tension between commercial speech rights and public-health efforts to curb overuse.

\section{Key Quotations}

\begin{quote}
``As to the particular consumer's interest in the free flow of commercial information, that interest may be as keen, if not keener by far, than his interest in the day's most urgent political debate. \dots\ So long as we preserve a predominantly free enterprise economy, the allocation of our resources in large measure will be made through numerous private economic decisions. It is a matter of public interest that those decisions, in the aggregate, be intelligent and well informed. To this end, the free flow of commercial information is indispensable.''
\end{quote}
\noindent --- Justice Blackmun, 425 U.S.\ at 763, 765.

\begin{quote}
``It is precisely this kind of choice, between the dangers of suppressing information, and the dangers of its misuse if it is freely available, that the First Amendment makes for us.''
\end{quote}
\noindent --- Justice Blackmun, 425 U.S.\ at 770.

\begin{quote}
``If there is a right to advertise, there is a reciprocal right to receive the advertising, and it may be asserted by these appellees.''
\end{quote}
\noindent --- Justice Blackmun, 425 U.S.\ at 757.

\end{document}
